\documentclass[10pt]{article}
\usepackage[T1]{fontenc}
\usepackage[utf8]{inputenc}
\usepackage[french]{babel}
\usepackage{amsmath,amssymb,mathrsfs,mathpazo}
\usepackage{geometry}
\geometry{a4paper,margin=2cm}
\usepackage{graphicx,array,tabularx,tabularray,colortbl,multirow,multicol,booktabs,enumitem,pifont,ulem,fancybox,diagbox,cancel}
\usepackage{pstricks,pstricks-add,pst-plot,pst-tree,pst-node,pst-text,pst-func,pst-math,pst-bspline,pst-3dplot,pst-eucl,pst-solides3d}
\let\pscalculate\undefined
\usepackage{pst-bezier}
\usepackage{tikz,pgfplots,tkz-tab}
\usetikzlibrary{arrows,arrows.meta,patterns,calc,positioning,trees,decorations.markings,shapes.geometric,shapes.misc,shapes.arrows,intersections,decorations.pathmorphing,decorations.pathreplacing,backgrounds,shadows,babel}
\pgfplotsset{compat=1.18}
\usepgfplotslibrary{fillbetween}
\usepackage[np]{numprint}
\usepackage{listings,makecell,pas-tableur,twoopt}
\let\barre\undefined
\usepackage[pstricks]{tablvar}
\definecolor{darkgreen}{RGB}{0,100,0}
\providecommand{\red}{\color{red}}
\providecommand{\blue}{\color{blue}}
\providecommand{\green}{\color{green}}
\providecommand{\black}{\color{black}}
\providecommand{\white}{\color{white}}
\providecommand{\gray}{\color{gray}}
\providecommand{\cyan}{\color{cyan}}
\providecommand{\magenta}{\color{magenta}}
\providecommand{\yellow}{\color{yellow}}
\providecommand{\violet}{\color{violet}}
\providecommand{\orange}{\color{orange}}
\providecommand{\euro}{\texteuro}
\providecommand{\EUR}{\texteuro}
\providecommand{\ieme}{\textsuperscript{e}}
\providecommand{\degres}{\ensuremath{^\circ}}
\providecommand{\degree}{\ensuremath{^\circ}}
\providecommand{\vv}[1]{\overrightarrow{#1}}
\providecommand{\Coord}[1]{#1}
\providecommand{\decosix}{}
\providecommand{\starredbullet}{\ensuremath{\bullet}}
\pagestyle{empty}
\setlength{\parindent}{0pt}
\usepackage{adjustbox}\DeclareUnicodeCharacter{2236}{:}\DeclareUnicodeCharacter{2190}{\ensuremath{\leftarrow}}
\providecommand{\R}{\mathbb{R}}
\providecommand{\N}{\mathbb{N}}
\providecommand{\Z}{\mathbb{Z}}
\providecommand{\Q}{\mathbb{Q}}
\providecommand{\D}{\mathbb{D}}
\providecommand{\C}{\mathbb{C}}
\providecommand{\e}{{\mathrm{e}}}
\providecommand{\vect}[1]{\overrightarrow{#1}}
\providecommand{\vectt}[1]{\overrightarrow{\mathrm{#1}}}
\providecommand{\vvt}[1]{\overrightarrow{\mathrm{#1}}}
\providecommand{\np}[2][]{{#2~#1}}
\providecommand{\vv}[1]{{\overrightarrow{#1}}}
\renewcommand{\d}{{\mathrm{d}}}
\providecommand{\degres}{{^\circ}}
\providecommand{\degre}{{^\circ}}
\providecommand{\blue}{\color{blue}}
\providecommand{\red}{\color{red}}
\providecommand{\enskip}{\quad}
\providecommand{\llbracket}{\lbrack\!\lbrack}
\providecommand{\rrbracket}{\rbrack\!\rbrack}
\providecommand{\euro}{{\text{€}}}
\providecommand{\no}{{\text{n°}}}
\providecommand{\interval}[2]{{[#1;#2]}}
\providecommand{\intervFO}[2]{{[#1;#2[}}
\providecommand{\dotfill}{\ldots}

\providecommand{\ds}{\displaystyle}
\providecommand{\textding}[1]{#1}
\providecommand{\Oij}{\(\mathrm{O};\vec{\imath},\vec{\jmath}\)}
\providecommand{\Oijk}{\(\mathrm{O};\vec{\imath},\vec{\jmath},\vec{k}\)}
\providecommand{\Ouv}{\(\mathrm{O};\vec{u},\vec{v}\)}
\providecommand{\intervFF}[2]{\ensuremath{[#1;#2]}}
\renewcommand{\euro}{\ifmmode\text{\texteuro}\else\texteuro\fi}
\DeclareUnicodeCharacter{00B0}{\ifmmode^{\circ}\else\textdegree\fi}
\DeclareUnicodeCharacter{20AC}{\ifmmode\text{\texteuro}\else\texteuro\fi}
% Énoncés : Éducation nationale. Transcription et corrigés : APMEP, Denis Vergès et contributeurs. Découpage et classement : CoopMaths, Sébastien Lozano, Éric Elter et contributeurs. Adaptation de la présentation : Maths Explorer.
% CC BY-SA : https://coopmaths.fr/www/about/
% https://coopmaths.fr/alea/static/bac/2022/tex/bac_2022_09_sujet1_ameriquesud_2_cor.tex
\begin{document}
\noindent\textit{APMEP / CoopMaths. CC BY-SA. Présentation adaptée par Maths Explorer.}\par\medskip
\medskip

Soit $\left(u_n\right)$ la suite définie par $u_0 = 4$ et, pour tout entier naturel $n$,\: $u_{n+1} = \dfrac15 u_n^2$.



\begin{enumerate}
\item 
	\begin{enumerate}
		\item $u_1 = \dfrac{1}{5} u_0^2 = \dfrac{16}{5}$ et $u_2 = \dfrac{1}{5} u_1^{2} = \dfrac{\left (\frac{16}{5}\right )^2}{5} = \dfrac{256}{125}$
		
		\item On complète la fonction ci-dessous qui renvoie la valeur du terme de rang $p$ de la suite $\left(u_n\right)$.

\begin{center}
\begin{tabular}{l}
\texttt{def suite\_u(p) :}\\
\quad \texttt{u = \blue 4}\\
\quad \texttt{for i in range(1,{\blue p + 1}) :}\\
\quad \qquad \texttt{u = \blue u*u/5}\\
\quad \texttt{return u}\\
\end{tabular}
\end{center}
	\end{enumerate}

\item 
	\begin{enumerate}
		\item 
Soit $\mathcal{P}_n$ la propriété $0 < u_n \leqslant 4$.

\begin{list}{\textbullet}{}
\item \textbf{Initialisation}

$u_0=4$ et $0<4\leqslant 4$ donc $\mathcal{P}_0$ est vraie.

\item \textbf{Hérédité}

On suppose $\mathcal{P}_n$ vraie, c'est-à-dire $0 < u_n \leqslant 4$ (hypothèse de récurrence).

$0 < u_n \leqslant 4$ donc
$0 < u_n^2 \leqslant 16$ donc
$0 < \dfrac{1}{5} u_n^2 \leqslant \dfrac{16}{5}$

Or $\dfrac{16}{5} =3,2 \leqslant 4$ donc $0 < \dfrac{1}{5} u_n^2 \leqslant 4$, c'est-à-dire $0 < u_{n+1} \leqslant 4$.

Donc $\mathcal{P}_{n+1}$ est vraie.

\item \textbf{Conclusion}

La propriété est vraie au rang 0, et elle est héréditaire. Donc, d'après le principe de récurrence, elle est vraie pour tout entier naturel.
\end{list}

On a donc démontré que pour tout entier naturel $n$, on a: $0 < u_n \leqslant 4$.

		\item 
$u_{n+1}-u_n = \dfrac{1}{5} u_n^2 - u_n = u_n \left (\dfrac{u_n}{5}-1\right )
= u_n \left (\dfrac{u_n-5}{5}\right )$		

$u_n\leqslant 4$ donc $u_n-5<0$; or $u_n>0$ donc $u_n\left (u_n-5\right )<0$ et donc $u_{n+1}-u_n<0$

On en conclut que la suite $(u_n)$ est décroissante.
		
		\item 
Pour tout $n$, $0<u_n$ donc la suite $(u_n)$ est minorée par 0.

La suite $(u_n)$ est décroissante et minorée donc, d'après le théorème de la convergence monotone, la suite $(u_n)$ est convergente.		
	\end{enumerate}
	
\item 
	\begin{enumerate}
		\item
\begin{list}{\textbullet}{}
\item $\ds\lim_{n\to +\infty} u_{n+1} = \ds\lim_{n\to +\infty} u_{n} =\ell$;
\item $\ds\lim_{n\to +\infty} \dfrac{1}{5}u_{n}^2  = \dfrac{1}{5}\ell^2$;
\item Pour tout $n$, on a: $u_{n+1}=\dfrac{1}{5}u_n^2$.
\end{list}		

On peut en conclure que $\ell = \dfrac{1}{5}\ell^2$.
		
		\item 
$\ell$ est solution de l'équation $x=\dfrac{1}{5}x^2$; on résout cette	équation.

$x=\dfrac{1}{5}x^2 \iff x\left (1-\dfrac{1}{5}x\right )=0
\iff x=0 \text{ ou } x=5$

Or, pour tout $n$, on a: $0<u_n\leqslant 4$ donc $\ell$ vérifie $0\leqslant \ell \leqslant 4$. La solution $\ell=5$ n'est donc pas valable donc $\ell=0$.
		
	\end{enumerate}
\item Pour tout entier naturel $n$, on pose $v_n = \ln \left(u_n\right)$ et $w_n = v_n - \ln (5)$.
	\begin{enumerate}
		\item 
$\begin{gathered}v_{n+1} = \ln \left (u_{n+1}\right ) = \ln \left (\dfrac{1}{5}u_n^2\right )
= \ln \left ( \dfrac{1}{5}\right ) + \ln \left ( u_n^2 \right )
= - \ln \left (5\right ) +2 \ln \left (u_n\right )
=2 \ln\left (u_n\right ) -\ln \left (5\right )\\
\phantom{v_{n+1}}
=2 v_n -\ln \left (5\right )\end{gathered}$		
		
		\item 
$w_n = v_n - \ln (5)$ donc $v_n=w_n+ \ln(5)$

$\begin{gathered}w_{n+1} =v_{n+1} - \ln(5) = 2 v_n -\ln(5) -\ln(5)
= 2\left (w_n +\ln(5)\mathstrut\right ) -2\ln(5)\\
\phantom{w_{n+1}}
=2 w_n +2\ln(5)-2\ln(5)=2w_n\end{gathered}$		
		
Donc la suite $\left(w_n\right)$ est géométrique de raison 2.
		
		\item 
$v_0=\ln\left (u_0\right )=\ln(4)$ donc $w_0=v_0-\ln(5)=\ln(4)-\ln(5) = \ln \left (\dfrac{4}{5}\right )$

La suite $(w_n)$ est géométrique de raison $q=2$ et de premier terme $w_0=\ln\left (\dfrac{4}{5}\right )$ donc, pour tout entier naturel $n$, on a:
$w_n=w_0\times q^n = \ln\left (\dfrac{4}{5}\right ) \times 2^n$.

Pour tout $n$, $v_n=w_n +\ln(5)$ donc $v_n = \ln\left (\dfrac{4}{5}\right )\times 2^n +\ln(5)$.		
			\end{enumerate}
			
\item 
$\ds\lim_{n\to +\infty} 2^n=+\infty$;
$\dfrac{4}{5}<1$ donc $\ln\left (\dfrac{4}{5}\right )<0$.

On en déduit que $\ds\lim_{n\to +\infty} \ln\left (\dfrac{4}{5}\right )\times 2^n=-\infty$, et donc que $\ds\lim_{n\to +\infty} v_n=-\infty$.

$v_n=\ln\left (u_n\right )$ donc $u_n=\e^{v_n}$

Or $\ds\lim_{x\to -\infty} \e^{x}=0$, et $\ds\lim_{n\to +\infty} v_n=-\infty$, donc $\ds\lim_{n\to +\infty} \e^{v_n}=0$.

On a donc redémontré que $\ds\lim_{n\to +\infty} u_n=0$.


\end{enumerate}

\vspace{0.5cm}

\end{document}
